% !TEX root = main_proposal.tex
% =====================================================================
%  PREAMBLE BERSAMA — dipakai oleh main_proposal.tex dan main_skripsi.tex
%  (\input dari kedua entry point supaya paket/style tidak dobel-tulis)
% =====================================================================

\documentclass[12pt,a4paper]{report}

% ── Bahasa & font ────────────────────────────────────────────────────
\usepackage[indonesian]{babel}   % otomatis: Bab / Gambar / Tabel / Daftar Isi
\usepackage{fontspec}            % wajib XeLaTeX/LuaLaTeX, bukan pdflatex
% Calibri terpasang di sistem ini, jadi dipakai langsung -- ini font yang
% diwajibkan Panduan Skripsi Filkom UB v3.0 (Lampiran A.3), bukan sekadar
% pengganti. Kalau compile di mesin lain yang TIDAK punya Calibri (mis.
% Overleaf, Linux, atau Mac tanpa MS Office), ganti "Calibri" -> "Carlito"
% pada baris \setmainfont di bawah -- itu pengganti metrik-kompatibel
% resmi untuk Calibri (dipakai LibreOffice sebagai substitusinya),
% sehingga tata letak/line-break tetap identik.
\setmainfont{Calibri}
% Catatan: paket sansmathfonts (font matematika sans-serif) tidak tersedia
% di lingkungan kompilasi ini dan dihapus. Simbol matematika pada equation
% akan memakai font matematika default (Latin Modern Math, serif tipis)
% sementara badan teks tetap Calibri — ini perilaku standar LaTeX dan tetap
% rapi dibaca. Jika ingin font matematika sans-serif penuh, pasang paket
% "sansmath" atau "unicode-math" secara manual di sistem Anda.
\usepackage{setspace}
\singlespacing                   % spasi 1 — Panduan Skripsi Filkom UB v3.0, Lampiran A.4

% ── Jarak baris daftar tabel dan daftar gambar ───────────────────────────────────────────────
\usepackage{etoolbox}
\makeatletter
\patchcmd{\@chapter}{\addtocontents{lot}{\protect\addvspace{10\p@}}}{}{}{}
\patchcmd{\@chapter}{\addtocontents{lof}{\protect\addvspace{10\p@}}}{}{}{}
\makeatother

% ── Tata letak halaman ───────────────────────────────────────────────
\usepackage[a4paper,left=4cm,right=3cm,top=3cm,bottom=3cm]{geometry}
\setlength{\parindent}{1.27cm}
\setlength{\parskip}{8pt plus 2pt minus 2pt}   % jarak vertikal antarparagraf, mengikuti template resmi

% ── Kerapatan daftar bernomor/berpoin ────────────────────────────────
% enumitem dipakai supaya jarak antarbutir enumerate/itemize tidak ikut
% melebar akibat \parskip di atas -- tanpa ini, tiap butir daftar akan
% mendapat jarak sebesar \parskip juga, jauh lebih lebar dari yang wajar.
\usepackage{enumitem}
\setlist{itemsep=2pt, parsep=0pt, topsep=4pt, partopsep=0pt}

% ── Tabel & gambar ───────────────────────────────────────────────────
\usepackage{graphicx}
\graphicspath{{assets/}}
\usepackage{booktabs}
\usepackage{array}
\usepackage{multirow}
\usepackage{longtable}
\usepackage{float}
\usepackage{caption}
% Semua caption: bold label, centered, posisi default (bawah untuk gambar, atas untuk tabel)
\captionsetup{font=normalsize, labelfont=bf, textfont=bf, justification=centering, singlelinecheck=false, labelsep=space}
\captionsetup[figure]{position=below}
\captionsetup[table]{position=above}

% ── Tipe float kustom (ganti nama caption sesuai kebutuhan) ──────────
% Contoh: \DeclareCaptionType{diagram}[Diagram][Daftar Diagram]
% Contoh: \DeclareCaptionType{bagan}[Bagan][Daftar Bagan]
% Setelah deklarasi, gunakan environment {diagram} seperti {figure}
%   \begin{diagram}[H]
%     \centering
%     \includegraphics[width=\textwidth]{nama_file}
%     \caption{Judul diagram}
%   \end{diagram}

% ── Tabel full-border, boleh sambung ke halaman berikutnya ───────────
% Iterasi 10 item 11: semua tabel pakai longtable (bukan table+tabular)
% supaya baris otomatis lanjut ke halaman berikutnya kalau kepanjangan,
% dengan header berulang di tiap halaman lanjutan. Garis vertikal
% (|c|c|) dan \hline di tiap baris dipertahankan di semua tabel --
% booktabs/adjustbox TIDAK dipakai lagi untuk tabel isi bab.
% Template:
%   \begin{longtable}{|p{4cm}|p{9cm}|}
%     \caption{Judul Tabel}
%     \label{tab:label} \\
%     \hline
%     \textbf{Kolom 1} & \textbf{Kolom 2} \\ \hline
%     \endfirsthead
%     \hline
%     \textbf{Kolom 1} & \textbf{Kolom 2} \\ \hline
%     \endhead
%     \hline
%     \endfoot
%     \hline
%     \endlastfoot
%     Isi baris & Isi baris \\ \hline
%   \end{longtable}

% ── Matematika ───────────────────────────────────────────────────────
\usepackage{amsmath}
\numberwithin{equation}{chapter}   % nomor persamaan per-bab: (3.1), (3.2)
\usepackage{amssymb}
\usepackage{siunitx}
\sisetup{output-decimal-marker={,}}

% ── Heading bab & subbab ─────────────────────────────────────────────
\usepackage{titlesec}
\usepackage{indentfirst}         % inden paragraf pertama juga (bukan cuma paragraf ke-2 dst),
                                  % standar konvensi Indonesia -- default LaTeX tidak begini
% Bab: "BAB X JUDUL" satu baris, centered, bold, normal size
% CATATAN: JANGAN pakai shape [block] di sini -- kombinasi [block] + \chapter
% di kelas report bikin \tableofcontents/\listoffigures/\listoftables render
% KOSONG (isi .toc/.lof/.lot tetap benar, cuma tidak ke-render ke halaman).
% Tanpa opsi shape (default titlesec), tampilan visualnya identik tapi TOC
% normal. Sudah diverifikasi lewat isolasi bug, jangan tambahkan [block] lagi.
% Ukuran font heading (Iterasi 10 item 4): bab 16pt, subbab 14pt, sub-subbab 12pt
\titleformat{\chapter}
  {\normalfont\fontsize{16}{20}\selectfont\bfseries\centering}
  {BAB \thechapter\space}
  {0pt}
  {\MakeUppercase}
\titlespacing*{\chapter}{0pt}{-40pt}{20pt}
% Subbab: "X.Y Judul" bold, left-aligned, jarak label-judul normal
\titleformat{\section}
  {\normalfont\fontsize{14}{17}\selectfont\bfseries}
  {\thesection}
  {1em}
  {}
\titleformat{\subsection}
  {\normalfont\fontsize{12}{15}\selectfont\bfseries}
  {\thesubsection}
  {1em}
  {}

% \titleformat{\chapter} above only reliably restyles the NUMBERED form.
% \chapter* (used by \tableofcontents/\listoffigures/\listoftables and by
% "Daftar Referensi") falls back to the class's own \@makeschapterhead,
% which is left-aligned/not-uppercase by default -- Iterasi 10 item 1
% flagged exactly this mismatch against "BAB 1 PENDAHULUAN"'s style.
% Redefine it directly so every unnumbered chapter heading matches: centered,
% same 16pt size, uppercase, same 20pt space-after as numbered chapters.
\makeatletter
\renewcommand{\@makeschapterhead}[1]{%
  {\parindent\z@ \centering \normalfont\fontsize{16}{20}\selectfont\bfseries
   \MakeUppercase{#1}\par\nobreak\vskip 20pt}%
}
\makeatother

% ── Daftar Isi / Daftar Gambar / Daftar Tabel ────────────────────────
\usepackage{tocloft}
% Tanpa bold di semua entri (termasuk level bab, yang default-nya bold)
\renewcommand{\cftchapfont}{\normalfont}
\renewcommand{\cftchappagefont}{\normalfont}
% Entri bab di daftar isi ikut menuliskan "BAB X", sama seperti heading aslinya
\renewcommand{\cftchappresnum}{BAB\space}
\setlength{\cftchapnumwidth}{3em}
% Titik-titik penghubung ke nomor halaman juga muncul di level bab
% (default kelas report: level bab tanpa titik-titik, cuma section ke bawah)
\renewcommand{\cftchapdotsep}{\cftdotsep}
% Titik-titik lebih rapat di semua level
\renewcommand{\cftdotsep}{1}
% Jarak vertikal antarentri dipersempit lagi (Iterasi 13 item 1 -- 2pt masih
% dirasa lebar setelah dicoba) -- 0pt di semua level: bab (Daftar Isi),
% gambar (Daftar Gambar), tabel (Daftar Tabel), dan section (baris "X.Y
% Judul" di dalam Daftar Isi, belum pernah disentuh sebelumnya)
\setlength{\cftbeforechapskip}{0pt}
\setlength{\cftbeforefigskip}{0pt}
\setlength{\cftbeforetabskip}{0pt}
\setlength{\cftbeforesecskip}{0pt}

% ── Mikrografi & hyperlink ───────────────────────────────────────────
\usepackage{microtype}
\setlength{\emergencystretch}{3em}
\usepackage{url}                 % \url{} bisa memecah baris di URL/DOI panjang
\PassOptionsToPackage{hyphens}{url}
\usepackage[hidelinks]{hyperref}
\hypersetup{
  pdftitle={Pengaruh Skema Pembagian Data Berbasis Film pada Dataset E-SERAVD terhadap Performa Klasifikasi Pengenalan Emosi Ucapan Berbahasa Indonesia},
  pdfauthor={Gaung Taqwa Indraswara},
}

% ── Warna ────────────────────────────────────────────────────────────
\usepackage{xcolor}
\definecolor{cProses}{HTML}{4A6D8C}    \definecolor{cProsesBg}{HTML}{DCE6F0}
\definecolor{cData}{HTML}{8C7B4A}      \definecolor{cDataBg}{HTML}{F0EBDC}
\definecolor{cKeputusan}{HTML}{8C6A4A} \definecolor{cKeputusanBg}{HTML}{F5E9D8}
\definecolor{cTerm}{HTML}{4A4A4A}      \definecolor{cTermBg}{HTML}{E4E4E4}
\definecolor{cGrup}{HTML}{9AA7B4}      \definecolor{cGrupBg}{HTML}{F2F5F8}
\definecolor{cOk}{HTML}{3E7D5A}        \definecolor{cOkBg}{HTML}{DFF0E6}
\definecolor{cBad}{HTML}{A4503C}       \definecolor{cBadBg}{HTML}{F2DED7}
\definecolor{kodeBg}{HTML}{F7F7F4}
\definecolor{kodeKomentar}{HTML}{5A7D5A}
\definecolor{kodeKunci}{HTML}{4A4A8C}
\definecolor{kodeString}{HTML}{8C5A3C}

% ── TikZ / pgfplots ──────────────────────────────────────────────────
\usepackage{tikz}
\usepackage{pgfplots}
\pgfplotsset{compat=1.18}
\usetikzlibrary{
  shapes.geometric, arrows.meta, positioning,
  calc, fit, backgrounds, decorations.pathreplacing
}

% node styles (flowchart monokrom)
\tikzset{
  term/.style={ellipse, draw=black, line width=.5pt,
    minimum width=26mm, minimum height=9mm,
    align=center, font=\footnotesize},
  proc/.style={rectangle, draw=black, line width=.5pt,
    minimum width=34mm, minimum height=11mm,
    inner xsep=3.4mm, align=center, font=\footnotesize,
    append after command={\pgfextra{
      \draw[line width=.5pt]
        ([xshift=1.6mm]\tikzlastnode.north west)
        -- ([xshift=1.6mm]\tikzlastnode.south west);
      \draw[line width=.5pt]
        ([xshift=-1.6mm]\tikzlastnode.north east)
        -- ([xshift=-1.6mm]\tikzlastnode.south east);}}},
  procplain/.style={rectangle, draw=black, line width=.5pt,
    minimum width=34mm, minimum height=11mm,
    align=center, font=\footnotesize},
  io/.style={trapezium, trapezium left angle=70, trapezium right angle=110,
    draw=black, line width=.5pt,
    minimum width=28mm, minimum height=9mm,
    align=center, inner sep=2pt, font=\footnotesize},
  dec/.style={diamond, aspect=2, draw=black, line width=.5pt,
    align=center, inner sep=1pt, font=\footnotesize},
  arr/.style={-{Stealth[length=2.2mm]}, draw=black, line width=.5pt},
  flbl/.style={font=\footnotesize, fill=white, inner sep=1pt},
}

% ── Listing kode ─────────────────────────────────────────────────────
\usepackage{listings}
\renewcommand{\lstlistingname}{Kode Sumber}
\lstdefinestyle{kode}{
  backgroundcolor=\color{kodeBg},
  basicstyle=\ttfamily\footnotesize,
  commentstyle=\color{kodeKomentar}\itshape,
  keywordstyle=\color{kodeKunci}\bfseries,
  stringstyle=\color{kodeString},
  numberstyle=\ttfamily\scriptsize\color{gray},
  numbers=left, numbersep=8pt, stepnumber=1,
  frame=single, rulecolor=\color{cGrup},
  breaklines=true, breakatwhitespace=false,
  showstringspaces=false, tabsize=2,
  captionpos=t, xleftmargin=18pt, framexleftmargin=14pt,
  literate={→}{$\rightarrow$}{1} {≤}{$\le$}{1} {≥}{$\ge$}{1}
           {×}{$\times$}{1} {—}{{-{}-}}1 {–}{{-}}1
}
\lstset{style=kode}

% ── Makro bantu ──────────────────────────────────────────────────────
\usepackage[export]{adjustbox}
\newcommand{\gambarbesar}[1]{%
  \adjustbox{max width=\textwidth, max totalheight=0.82\textheight, center}{#1}}

% ── Jenis dokumen (diisi beda oleh main_proposal.tex / main_skripsi.tex) ──
% \doctypelabel: label besar di cover ("PROPOSAL SKRIPSI" atau "SKRIPSI").
% \doctypesubtitle + \ifdoctypesubtitle: baris kecil di bawahnya, cuma
% dipakai skripsi penuh -- boolean-nya yang menentukan baris itu dicetak
% atau tidak (bukan sekadar dikosongkan), supaya versi proposal tidak
% menyisakan jarak kosong dari baris yang seharusnya tidak ada.
\newcommand{\doctypelabel}{}
\newcommand{\doctypesubtitle}{}
\newif\ifdoctypesubtitle
\doctypesubtitlefalse

% ── Metadata dokumen ─────────────────────────────────────────────────
\title{Pengaruh Skema Pembagian Data Berbasis Film pada Dataset E-SERAVD
  terhadap Performa Klasifikasi Pengenalan Emosi Ucapan Berbahasa Indonesia}
\author{Gaung Taqwa Indraswara}
\date{\the\year}
